\chapter{为什么人工装入需要变成机器自己的工作}

\begin{keyidea}
操作人员能在复位期间把一项程序、输入和初始栈写入 RAM；释放复位
后，启动程序建立寄存器并移交处理器。程序可以计算并合作返回，但更换程序仍要停机并借助
外部写入工具。本章增加运行中的字节接收与装入事务，不改变单任务执行模型。
\end{keyidea}

现在只改变一个条件：机器已经装入机箱，运行期间不能再由操作人员逐地址改写 RAM。外部
设备仍然能够按顺序送来字节，但机器必须自己接收、验证并放置这些字节。为此，本章才增加
串行控制器和装入程序；它们不负责同时保留两项运行现场，也不能在程序失控时收回处理器。

\section{从人工写入转向运行中的字节接收}

上一章的操作人员可以在复位期间逐地址写 RAM。机器一旦离开实验台，这种方式就出现了
具体成本：每次更换程序都要停机、连接调试工具、重新写入多个区域，并由人保证最后一个
字节已经完成。机器自身无法重复这套过程。

一种选择是把永远不变的程序直接固化在启动存储中。这样不再需要每次装入，但更换程序要
重新写非易失器件；输入参数和可写栈仍要另行准备。本章需要的是“保持固定启动程序不变，
运行期间接收另一份可替换程序”，所以必须给运行中的处理器增加一个能够取得外部字节的
接口。

\subsection{为什么不能收到一条指令就立刻执行}

外部连接按时间送来字节，处理器执行程序却会分支回到先前地址。若机器收到四个指令字节就
立即执行，循环第二次取指时，先前字节已经从顺序数据流中消失。除非外部端同时承担一片可
随机访问的指令存储，否则流式执行无法支持本书的求和循环。

另一个问题是完整性。前几条指令到达后便开始执行，后续传输若中断或校验失败，已经执行的
指令可能改写 RAM。此时不能再把整次接收当作“从未发生”。因此，基本方案先把全部程序
写入 RAM，验证通过以后才改变处理器入口。

\begin{longtable}{|L{0.24\linewidth}|L{0.28\linewidth}|L{0.34\linewidth}|}
\hline
\rowcolor{BookTableHead}
方案 & 省去的步骤 & 失去的能力或边界 \\ \hline
程序固化在启动存储 & 每次传输与装入 & 不能方便替换程序 \\ \hline
收到一条执行一条 & 完整 RAM 映像 & 不能可靠分支回旧指令 \\ \hline
接收与执行重叠 & 等待完整传输的时间 & 不能在产生效果前完成整体验证 \\ \hline
先接收、验证、再移交 & 无 & 需要接收状态、RAM 区域和明确提交点 \\ \hline
\end{longtable}

最后一行没有宣称适用于所有机器。它只是给本章提供一条可检查的基本路径：外部字节先成为
完整映像，处理器随后一次性从启动程序切到新入口。串行控制器只负责把线上位流整理成字节；
映像格式、范围检查和最终移交仍由启动程序负责。

\section{让外部任务进入机器：串行控制器}

可以把求和任务永久写进启动存储，但那样每换一次任务都要重新烧写固件。本章为机器增加
一个最小串行控制器。外部发送端在一根接收线上按时间发送位，控制器把位组合成字节，再把
字节暴露为处理器可以寻址的状态。发送方向用于固件报告接受、拒绝和最终结果。

\subsection{新器件接在哪里}

串行控制器有两类连接。第一类是机器外部的收发信号；第二类是接入现有地址互连的配置与
数据端口。处理器不直接观察线上的每一位，而是通过内存映射寄存器读取已经组装好的字节。

\begin{center}
\begin{tikzpicture}[node distance=10mm and 9mm]
\node[stage,minimum width=2.6cm] (host) {外部发送端\\任务字节流};
\node[stage,minimum width=3.0cm,right=of host] (uart) {串行控制器\\状态与收发槽};
\node[stage,minimum width=2.8cm,right=of uart] (bus) {地址互连\\设备范围};
\node[stage,minimum width=2.6cm,right=of bus] (cpu) {处理器\\轮询搬运};
\draw[flow] ([yshift=3mm]host.east) -- node[above,font=\footnotesize]{接收线} ([yshift=3mm]uart.west);
\draw[flow] ([yshift=-3mm]uart.west) -- node[below,font=\footnotesize]{发送线} ([yshift=-3mm]host.east);
\draw[flow] ([yshift=3mm]cpu.west) -- node[above,font=\footnotesize]{读写请求} ([yshift=3mm]bus.east);
\draw[flow] ([yshift=-3mm]bus.east) -- node[below,font=\footnotesize]{响应} ([yshift=-3mm]cpu.west);
\draw[flow] ([yshift=3mm]bus.west) -- node[above,font=\footnotesize]{寄存器请求} ([yshift=3mm]uart.east);
\draw[flow] ([yshift=-3mm]uart.east) -- node[below,font=\footnotesize]{寄存器响应} ([yshift=-3mm]bus.west);
\end{tikzpicture}
\end{center}
\diagramnote{串行线与处理器请求不是同一种事件。控制器先在外部时间域接收位，再把完整字节
保存在内部槽中；处理器随后通过互连读取这个槽。}

\subsection{串行控制器内部保存什么}

\textbf{硬件模型。} 接收移位寄存器按串行时序收集位，形成一个字节后写入
\code{rx\_data} 并置 \code{rx\_ready=1}。处理器读取 \code{RXDATA} 后清除
\code{rx\_ready}。若旧字节尚未读取而新字节到达，控制器置
\code{rx\_overrun=1}。发送侧在 \code{tx\_ready=1} 时接受一个字节，移出期间清零
\code{tx\_ready}，发送完毕后重新置位。

\textbf{软件模型。} 地址互连为控制器分配三个 32 位寄存器：

\begin{longtable}{|L{0.28\linewidth}|L{0.18\linewidth}|L{0.42\linewidth}|}
\hline
\rowcolor{BookTableHead}
地址与名称 & 访问方式 & 软件可见语义 \\ \hline
\code{0x40000000 STATUS} & 只读 & bit0 为 \code{RX\_READY}，bit1 为
\code{RX\_OVERRUN}，bit2 为 \code{TX\_READY} \\ \hline
\code{0x40000004 RXDATA} & 只读 & 返回接收字节，并消费当前接收槽 \\ \hline
\code{0x40000008 TXDATA} & 只写 & 在 \code{TX\_READY=1} 时提交一个待发送字节 \\ \hline
\end{longtable}

固件读取字节时必须先检查状态：

\begin{lstlisting}
receive_byte:
  status = read32(0x40000000)
  if status has RX_OVERRUN: reject current transfer
  if status lacks RX_READY: repeat
  return read8(0x40000004)
\end{lstlisting}

\textbf{抽象。} 上层软件得到一个按到达次序读取、按提交次序发送的字节通道。它不保证
任务边界、长度、校验或重传；这些规则要由固件和外部发送端建立。它也不保证处理器在做
其他计算时仍能及时取走字节，单槽接收的速率上限由轮询速度决定。

\subsection{为什么不能只约定“发完为止”}

串行字节本身没有结束标记。若发送端静默，固件无法区分“任务已经发完”和“下一个字节
还没到”。因此，传输必须先提供长度。长度又可能因误码变成巨大数值，固件还需要固定标识
和校验，才能避免把随机字节当作可执行映像。

本章使用 32 字节的教学头部：

\begin{longtable}{|L{0.18\linewidth}|L{0.24\linewidth}|L{0.46\linewidth}|}
\hline
\rowcolor{BookTableHead}
偏移 & 字段 & 为什么必须存在 \\ \hline
\code{+0x00} & \code{magic} & 区分任务传输和随机输入，本章固定为 \code{MTK1} \\ \hline
\code{+0x04} & \code{header\_bytes} & 使接收者知道载荷从哪里开始 \\ \hline
\code{+0x08} & \code{image\_bytes} & 决定接收终点并用于范围检查 \\ \hline
\code{+0x0c} & \code{load\_address} & 指定载荷写入 RAM 的起始位置 \\ \hline
\code{+0x10} & \code{entry\_offset} & 在载荷内部指出第一条任务指令 \\ \hline
\code{+0x14} & \code{zero\_bytes} & 指定载荷后需清零的可写区域 \\ \hline
\code{+0x18} & \code{stack\_bytes} & 说明任务要求的最小栈空间 \\ \hline
\code{+0x1c} & \code{checksum} & 检测头部之外的载荷是否损坏 \\ \hline
\end{longtable}

这只是固件与发送端的装入协议，不是文件系统中的文件格式。它没有文件名、目录、权限、
时间戳或持久对象身份。给它增加这些名称并不会让机器自动获得文件系统。

\section{装入不是复制一遍就结束}

固件收到头部后，不能立即按其中地址写 RAM。一个恶意或损坏的
\code{image\_bytes} 可能在加法中溢出，使表面上的终点变成较小地址；载荷可能覆盖固件
工作区或任务栈；入口偏移也可能指到载荷之外。所有会影响写入范围和最终 PC 的字段都必须
先验证。

\subsection{先确定本次内存所有权}

本章采用以下 RAM 布局：

\begin{longtable}{|L{0.29\linewidth}|L{0.27\linewidth}|L{0.32\linewidth}|}
\hline
\rowcolor{BookTableHead}
物理范围 & 启动期间的所有者 & 用途 \\ \hline
\code{0x00100000--0x0017ffff} & 待装入任务 & 映像、零填充区和任务数据 \\ \hline
\code{0x00180000--0x001effff} & 待装入任务 & 向下增长的任务栈候选区 \\ \hline
\code{0x001f0000--0x001fffff} & 固件 & 接收状态、校验状态和固件栈 \\ \hline
\end{longtable}

“所有者”此时只是固件遵守的软件约定，并非硬件权限。固件在移交前不会主动把任务载荷写
进自己的工作区；任务开始后却仍可用普通存储指令覆盖那里。正是这种缺乏强制边界的事实，
决定了本章还不能称为有了操作系统保护。

\subsection{范围检查必须考虑整数溢出}

设
\[
L=\texttt{load\_address},\qquad
N=\texttt{image\_bytes},\qquad
Z=\texttt{zero\_bytes}.
\]
若直接计算 \(L+N+Z\)，32 位加法可能回绕。稳妥检查先确认每个长度不超过允许区间，再用
减法形式判断：
\[
N \le U-L,\qquad
Z \le U-(L+N),
\]
其中 \(U\) 是任务载入区的开区间上界 \code{0x00180000}。在使用 \(U-L\) 前，还要确认
\(L<U\) 且 \(L\) 不低于 \code{0x00100000}。

入口地址
\[
E=L+\texttt{entry\_offset}
\]
必须落在已接收的映像字节中，并满足 4 字节指令对齐。栈大小必须能在候选栈区内安排，且
栈区不能与载荷及零填充区重叠。

\begin{lstlisting}
validate_header(h):
  require h.magic == MTK1
  require h.header_bytes == 32
  require TASK_LOAD_LOW <= h.load_address < TASK_LOAD_HIGH
  require h.image_bytes <= TASK_LOAD_HIGH - h.load_address
  image_end = h.load_address + h.image_bytes
  require h.zero_bytes <= TASK_LOAD_HIGH - image_end
  require h.entry_offset < h.image_bytes
  require (h.load_address + h.entry_offset) is 4-byte aligned
  require MIN_STACK <= h.stack_bytes <= STACK_REGION_SIZE
\end{lstlisting}

验证顺序是协议的一部分。只有前一项证明相应运算安全，后一项才可以使用计算结果。把所有
条件压成一个长布尔表达式，会掩盖哪一步可能先发生溢出。

\subsection{逐字节接收时谁拥有哪些状态}

头部通过后，固件建立一个装入记录：

\begin{lstlisting}
load_record:
  destination_start
  destination_next
  payload_remaining
  checksum_so_far
  expected_checksum
  entry_address
  zero_start
  zero_bytes
  stack_low
  stack_top
  state = RECEIVING
\end{lstlisting}

\code{destination\_next} 指出下一个字节写到哪里；
\code{payload\_remaining} 防止接收循环依赖外部静默；
\code{checksum\_so\_far} 只覆盖已经接受的载荷。每消费一个 \code{RXDATA} 字节，固件按
固定顺序更新这些字段：

\begin{enumerate}
\item 从控制器消费一个字节；
\item 把字节写到 \code{destination\_next}；
\item 将该字节并入校验状态；
\item 递增 \code{destination\_next}；
\item 递减 \code{payload\_remaining}。
\end{enumerate}

若在第 2 步之后复位，RAM 中可能留下部分载荷，但复位会重新从启动存储执行，装入记录也
会重建。固件不能仅凭 RAM 中“看起来像代码”的字节跳转，必须重新完成协议验证。

\subsection{接收完成、验证完成与可执行不是一个时刻}

最后一个字节写入 RAM 时，只能说明传输长度已经满足。固件还要比较最终校验值、清零
\code{zero\_bytes} 指定的区域、建立输入数组和初始栈。为避免把半成品当作任务，装入
记录按以下状态推进：

\begin{center}
\begin{tikzpicture}[node distance=10mm and 12mm]
\node[stage,minimum width=2.8cm] (header) {等待头部};
\node[stage,minimum width=2.8cm,right=of header] (recv) {接收载荷};
\node[stage,minimum width=2.8cm,right=of recv] (verify) {校验与清零};
\node[stage,minimum width=2.8cm,right=of verify] (ready) {可移交};
\node[stage,minimum width=11.2cm,below=13mm of recv,fill=BookWarm] (reject)
  {任一头部、长度、校验或布局检查失败：拒绝当前传输并报告原因};
\draw[flow] (header) -- node[above,font=\scriptsize]{头部合法} (recv);
\draw[flow] (recv) -- node[above,font=\scriptsize]{长度满足} (verify);
\draw[flow] (verify) -- node[above,font=\scriptsize]{校验通过} (ready);
\draw[->,>=Latex,thick,draw=BookMuted] (header.south) -- (header.south |- reject.north);
\draw[->,>=Latex,thick,draw=BookMuted] (recv.south) -- (recv.south |- reject.north);
\draw[->,>=Latex,thick,draw=BookMuted] (verify.south) -- (verify.south |- reject.north);
\end{tikzpicture}
\end{center}
\diagramnote{“载荷到齐”只关闭接收阶段；只有校验、零填充和布局建立全部成功，装入记录才
进入“可移交”。任何失败都不能改变 PC 到任务入口。}

本章的提交点不是写入最后一个字节，而是稍后把处理器 PC 改为已经验证的入口地址。在此
之前，即使 RAM 中已有完整任务，处理器仍执行固件。

\subsection{用具体数字走完一次装入}

求和任务头部给出：

\begin{lstlisting}
magic          = MTK1
header_bytes   = 32
image_bytes    = 0x00000300
load_address   = 0x00100000
entry_offset   = 0x00000000
zero_bytes     = 0x00000100
stack_bytes    = 0x00004000
checksum       = expected value
\end{lstlisting}

载荷占用 \code{0x00100000--0x001002ff}，零填充区占用
\code{0x00100300--0x001003ff}，都低于任务载入区上界。任务栈使用
\code{0x001ec000--0x001effff}，与映像不重叠。入口地址就是
\code{0x00100000}。

输入数组是任务数据的一部分，位于 \code{0x00102000}，不在上述 0x300 字节载荷范围内
会产生矛盾。为使布局一致，本次实际映像应覆盖到输入区之后，因此把
\code{image\_bytes} 修正为 \code{0x00002100}。这项复算刻意展示了字段之间必须相互
核对：若只逐项抄表，装入器会接受一个宣称包含输入数据、实际长度却到不了输入地址的映像。

修正后的范围为：

\begin{longtable}{|L{0.32\linewidth}|L{0.22\linewidth}|L{0.32\linewidth}|}
\hline
\rowcolor{BookTableHead}
范围 & 装入后的内容 & 来源 \\ \hline
\code{0x00100000--0x0010001f} & 求和核心指令 & 载荷 \\ \hline
\code{0x00100020--0x00101fff} & 其他任务代码与常量 & 载荷 \\ \hline
\code{0x00102000--0x0010200f} & 四个输入整数 & 载荷 \\ \hline
\code{0x00102020--0x00102023} & 结果槽，初值为零 & 载荷 \\ \hline
\code{0x00102100--0x001021ff} & 未初始化区 & 固件清零 \\ \hline
\code{0x001ec000--0x001effff} & 任务栈 & 固件保留并建立初始内容 \\ \hline
\end{longtable}

外部发送端不应发送含糊的“差不多大小”。装入协议中的长度必须覆盖实际传输的每个字节，
固件也只能在校验过的范围内解释入口和数据。

\section{从固件现场建立任务现场}

任务字节可执行后，处理器仍保存固件的 PC、SP 和寄存器。直接把 PC 改为任务入口会让任务
继承任意寄存器值，并继续使用固件栈。任务第一次访问参数或执行 \code{RET} 时就会得到
错误结果。固件必须先构造一份明确的初始现场。

\subsection{初始寄存器约定}

本次启动约定如下：

\begin{longtable}{|L{0.18\linewidth}|L{0.26\linewidth}|L{0.44\linewidth}|}
\hline
\rowcolor{BookTableHead}
状态 & 移交值 & 原因 \\ \hline
\code{r0} & \code{0x00102000} & 输入数组首地址 \\ \hline
\code{r1} & \code{4} & 数组元素个数 \\ \hline
\code{r2} & \code{0x00102020} & 结果槽地址 \\ \hline
\code{r3--r7} & \code{0} & 消除对固件残留值的依赖 \\ \hline
\code{SP} & \code{0x001efffc} & 指向预置的固件返回地址 \\ \hline
\code{PC} & \code{0x00100000} & 最后一步才写入的任务入口 \\ \hline
\end{longtable}

SP 不是直接设为 \code{0x001f0000}。固件先在
\code{0x001efffc--0x001effff} 写入 32 位返回地址 \code{0x00000300}，再让 SP 指向这
个字。任务最终执行 \code{RET} 时，处理器读取 \code{MEM32[SP]} 作为新 PC，并把 SP 增加
4，于是控制流回到固件返回入口。

\begin{center}
\begin{tikzpicture}[x=1cm,y=0.72cm]
\draw[BookRule,thick] (0,0) rectangle (8,4.1);
\draw[BookRule] (0,1.1) -- (8,1.1);
\draw[BookRule] (0,2.2) -- (8,2.2);
\draw[BookRule] (0,3.3) -- (8,3.3);
\node[anchor=west] at (0.2,3.7) {\code{0x001f0000}：栈区上界，不属于栈中有效字};
\node[anchor=west] at (0.2,2.75) {\code{0x001efffc}：\code{0x00000300}，固件返回地址};
\node[anchor=west] at (0.2,1.65) {\code{0x001efff8}：后续调用可使用};
\node[anchor=west] at (0.2,0.55) {\code{0x001ec000}：栈区下界};
\draw[flow] (9.4,2.75) -- (8.05,2.75);
\node[anchor=west] at (9.5,2.75) {初始 SP};
\draw[->,>=Latex,BookMuted,thick] (8.8,1.7) -- (8.8,0.5);
\node[anchor=west,font=\footnotesize] at (9.0,1.1) {压栈方向};
\end{tikzpicture}
\end{center}
\diagramnote{向下增长表示压入新内容时 SP 减小。栈区上界本身是开区间；预置返回地址占用
上界下方最后 4 字节。}

\subsection{为什么入口地址必须最后提交}

建立初始现场的安全顺序是：

\begin{enumerate}
\item 确认装入记录处于“可移交”；
\item 停止接收本次映像，避免后续字节继续改写任务区；
\item 写入输入数据、结果初值和零填充区；
\item 建立任务栈并写入固件返回地址；
\item 设置通用寄存器和 SP；
\item 最后把 PC 改为任务入口。
\end{enumerate}

前五步仍在固件指令流中执行。第六步一旦提交，下一次取指来自任务区域，固件就不能继续
执行“还有一步初始化”。因此 PC 的改变是控制权移交的提交边界。

\begin{center}
\begin{tikzpicture}[node distance=10mm and 8mm]
\node[stage,minimum width=2.5cm] (loaded) {映像已验证};
\node[stage,minimum width=2.6cm,right=of loaded] (memory) {数据与栈已建立};
\node[stage,minimum width=2.6cm,right=of memory] (regs) {参数寄存器已写入};
\node[stage,minimum width=2.5cm,right=of regs] (pc) {PC 提交入口};
\node[stage,minimum width=3.1cm,below=14mm of pc] (task) {任务第一条取指};
\draw[flow] (loaded) -- (memory);
\draw[flow] (memory) -- (regs);
\draw[flow] (regs) -- node[above,font=\scriptsize]{移交点} (pc);
\draw[flow] (pc) -- (task);
\end{tikzpicture}
\end{center}
\diagramnote{“字节已经在 RAM 中”和“任务已经运行”之间隔着初始现场建立与 PC 提交。
只有最右侧事件改变取指来源。}

\subsection{移交前后的处理器快照}

移交前，固件可能有如下状态：

\begin{lstlisting}
PC = 0x000002d8       SP = 0x001ffec0
r0 = load_record_ptr r1 = checksum_so_far
r2 = 0x40000000      r3 = temporary value
\end{lstlisting}

移交完成后的第一条任务指令开始前，状态必须变为：

\begin{lstlisting}
PC = 0x00100000       SP = 0x001efffc
r0 = 0x00102000       r1 = 4
r2 = 0x00102020       r3 = 0
r4 = 0                r5 = 0
r6 = 0                r7 = 0
\end{lstlisting}

这两个快照说明“切换”不能只描述为跳转。至少有 PC、SP、参数寄存器和清理后的临时寄存器
发生变化。当前固件不需要保存自己的全部现场，因为任务返回后固件进入一个固定返回入口，
而不是回到 \code{0x000002d8} 继续原函数。下一章需要暂停并继续不同任务时，保存现场的
要求会更严格。

\section{把装入器写成可审计的状态机}

“固件接收并装入任务”仍然包含多个可能暂停或失败的阶段。若所有局部变量只散落在寄存器
中，错误入口无法判断已经接收多少字节，也无法给出稳定诊断。固件因此在自己的 RAM 工作区
保存一份装入记录。它不是操作系统任务对象，只代表当前这一笔传输。

\subsection{字段、写入者和有效期}

\begin{longtable}{|L{0.24\linewidth}|L{0.25\linewidth}|L{0.37\linewidth}|}
\hline
字段 & 谁在何时写入 & 后续读取目的 \\ \hline
\code{state} & 阶段转换代码 & 限制当前允许的操作 \\ \hline
\code{destination\_next} & 每接受一个字节后递增 & 决定下一笔 RAM 写地址 \\ \hline
\code{payload\_remaining} & 头部建立，逐字节递减 & 判断载荷何时到齐 \\ \hline
\code{checksum\_so\_far} & 每个已接受字节更新 & 与头部期望值比较 \\ \hline
\code{entry\_address} & 头部验证后计算一次 & 移交时写入 PC \\ \hline
\code{stack\_low/top} & 布局验证后计算一次 & 清零栈并设置 SP \\ \hline
\code{error\_code} & 首个失败检测点 & 报告哪条契约被破坏 \\ \hline
\end{longtable}

装入记录从收到头部开始存在，任务入口提交前仍由固件独占。任务运行后，记录不再描述一个
“正在接收”的事务；固定返回入口可以把它改为 \code{FINISHED}，或直接重置为
\code{WAITING\_HEADER}。若任务永不返回，记录也不会自行变化。

\subsection{允许的状态转换}

\begin{longtable}{|L{0.19\linewidth}|L{0.25\linewidth}|L{0.29\linewidth}|L{0.17\linewidth}|}
\hline
当前状态 & 输入事件 & 必须完成的动作 & 后继状态 \\ \hline
\code{WAITING} & 32 字节头到齐 & 解码并验证所有范围 & \code{RECEIVING} 或
\code{REJECTED} \\ \hline
\code{RECEIVING} & 一个载荷字节到达 & 写 RAM、更新校验和与计数 & 保持或
\code{VERIFYING} \\ \hline
\code{RECEIVING} & 控制器溢出 & 记录传输错误 & \code{REJECTED} \\ \hline
\code{VERIFYING} & 校验匹配 & 清零、建栈、建立参数 & \code{READY} \\ \hline
\code{VERIFYING} & 校验不匹配 & 记录内容错误 & \code{REJECTED} \\ \hline
\code{READY} & 固件决定运行 & 设置现场并提交 PC & \code{RUNNING} \\ \hline
\code{RUNNING} & 固定返回入口获得控制 & 读取并发送结果 & \code{FINISHED} \\ \hline
\end{longtable}

状态表排除了若干非法动作：\code{WAITING} 状态不能消费载荷；\code{REJECTED} 状态不能
跳任务入口；\code{RUNNING} 状态下固件也没有机会继续修改记录。把这些限制写入转换函数，
比在各处检查零散布尔量更容易审计。

\subsection{头部在串行线上也是字节}

以修正后的映像为例，头部前 28 个字节可表示为：

\begin{lstlisting}
offset  bytes          decoded value
00      4d 54 4b 31    magic "MTK1"
04      20 00 00 00    header_bytes = 0x20
08      00 21 00 00    image_bytes  = 0x2100
0c      00 00 10 00    load_address = 0x00100000
10      00 00 00 00    entry_offset = 0
14      00 01 00 00    zero_bytes   = 0x100
18      00 40 00 00    stack_bytes  = 0x4000
1c      .. .. .. ..    checksum
\end{lstlisting}

\code{load\_address} 的四个字节为 \code{00 00 10 00}。按低有效字节在前组合后才是
\code{0x00100000}。若固件按相反字节序解码，会得到 \code{0x00001000}，落入未映射
区域。传输协议必须规定字节序，不能依赖发送端和处理器“碰巧相同”。

\subsection{校验和覆盖什么，不能证明什么}

本章采用一个教学化的 32 位逐字节滚动校验：

\begin{lstlisting}
checksum = 0
for each payload byte b:
  checksum = rotate_left(checksum, 5)
  checksum = checksum XOR b
\end{lstlisting}

发送端对全部载荷计算结果并放入头部，固件按实际收到的次序重新计算。它能够高概率检测
常见误码、丢字节和次序变化，但不是密码学签名。知道算法的人可以构造另一个产生相同校验
值的载荷；校验通过也不能证明代码可信、没有越界指令或实现了正确算法。

这里有三种不同判断：

\begin{enumerate}
\item 格式有效：头部字段可安全解释，范围没有越界；
\item 传输一致：载荷长度和校验符合发送端给出的值；
\item 程序可信：代码来源和行为满足安全策略。
\end{enumerate}

本章只完成前两项。当前机器没有签名验证、权限或沙箱，不能从“校验通过”推出“任务可以
任意访问机器而没有风险”。

\subsection{用失败用例检验头部验证顺序}

给装入器一组具体头部，比只读正常伪代码更容易发现边界错误。

\subsection{长度恰好到达上界}

若 \code{load\_address=0x0017ff00}、\code{image\_bytes=0x100}、
\code{zero\_bytes=0}，映像终点恰好为开区间上界 \code{0x00180000}，可以接受。最后一个
有效字节地址为 \code{0x0017ffff}。

若 \code{image\_bytes=0x101}，最后一个字节会落在为栈候选区保留的
\code{0x00180000}，必须拒绝。用开区间终点计算可以避免“最后一个地址加一”再次溢出。

\subsection{32 位回绕伪装成小终点}

设攻击头部为：

\begin{lstlisting}
load_address = 0xfffffff0
image_bytes  = 0x00000040
\end{lstlisting}

32 位直接相加得到 \code{0x00000030}。若装入器只检查“终点小于 RAM 上界”，会错误接受。
正确顺序先检查起点是否位于任务区；起点检查已经失败，根本不执行可能回绕的终点判断。

\subsection{入口位于零填充区}

若 \code{entry\_offset=image\_bytes+4}，入口落在固件清零的区域。虽然地址属于任务 RAM，
那里并不是发送端提供且校验过的指令，所以必须拒绝。入口条件是
\[
0\le\texttt{entry\_offset}<\texttt{image\_bytes},
\]
而不是小于 \(\texttt{image\_bytes}+\texttt{zero\_bytes}\)。

\subsection{入口未对齐}

\code{entry\_offset=2} 使 PC 成为 \code{0x00100002}。该地址位于载荷内部，却不满足
教学处理器的 4 字节取指对齐，必须在移交前拒绝。若等到处理器取指才发现，错误仍可检测，
但诊断会从“非法映像入口”退化成“处理器对齐故障”，丢失了协议层原因。

\subsection{栈与映像不重叠仍可能过小}

只检查栈区不重叠还不够。本章至少要预置 4 字节返回地址，因此
\code{stack\_bytes<4} 必须拒绝。实际约定采用更大的
\code{MIN\_STACK}，为任务的函数调用保留空间。最小值是启动协议的一部分，不由 RAM 自动
推断。

\begin{problemframe}
\textbf{复算点。} 对
\code{load\_address=0x00170000}、\code{image\_bytes=0x10000}，
映像终点为 \code{0x00180000}，仍合法；若再要求
\code{zero\_bytes=1}，则第一字节零区会进入栈候选范围，整个头部必须拒绝。
\end{problemframe}

\subsection{串行速度如何限制轮询装入器}

采用常见的 8N1 串行帧时，一个数据字节在线上占 1 个起始位、8 个数据位和 1 个停止位，
合计 10 位。若传输速率为 \(115200\) 位每秒，一个字节约每
\[
\frac{10}{115200}\ \mathrm{s}\approx86.8\ \mathrm{\mu s}
\]
到达一次。

假设处理器频率为 \(10\,\mathrm{MHz}\)，两个字节到达之间约有 868 个处理器周期。固件
读取状态、消费字节、写 RAM 和更新校验必须在这段预算内完成，或让发送端等待。若最坏路径
需要 950 周期，平均速度看似接近，单槽控制器仍会周期性溢出。

\subsection{平均成本不足以证明不会溢出}

轮询循环平时可能只需 200 周期，但遇到错误报告、长 RAM 等待或其他固件工作时延迟会变大。
不会溢出的条件约束的是两次消费之间的最长间隔，而不是平均每字节成本：
\[
T_{\mathrm{consume,max}} < T_{\mathrm{arrival}}.
\]
如果无法满足，可以选择让发送端依据硬件流量控制暂停、增加接收 FIFO，或以后引入中断和
DMA。每种方案都会新增状态和失败模式；当前样机选择最简单的发送端应答协议。

\subsection{逐块应答建立背压}

外部端每发送 64 字节后等待固件返回一个确认字节。固件只有在这 64 字节全部写入 RAM 且
未发生溢出时才确认。若返回拒绝码，外部端从整个映像开始重传。本章不实现分块恢复，因为
那需要块编号、重复检测和部分校验等新状态。

这个协议把发送速率约束反馈给发送端，称为背压。背压不提高控制器内部容量，却阻止外部端
无限快地提交数据。即便没有操作系统，异步生产者与有限缓冲之间也已经存在同样的资源问题。

\subsection{一次成功移交的详细不变量}

在提交任务 PC 之前，固件逐项确认：

\begin{enumerate}
\item 装入记录状态为 \code{READY}；
\item \code{destination\_next=load\_address+image\_bytes}；
\item \code{payload\_remaining=0}；
\item 实际校验值等于头部校验值；
\item 入口位于已校验映像且满足对齐；
\item 零填充区已经写零；
\item 输入数组和结果槽位于任务数据范围；
\item 栈范围与映像、固件工作区均不重叠；
\item \code{MEM32[0x001efffc]=0x00000300}；
\item 通用寄存器和 SP 已按启动契约设置。
\end{enumerate}

这些条件称为移交不变量：只要固件准备改变 PC，它们就必须同时成立。若检查分散在不同函数
中，最终移交函数仍应以断言或显式返回值确认它们，而不能依赖“前面大概检查过”。

\begin{center}
\begin{tikzpicture}[node distance=8mm and 12mm]
\node[stage,minimum width=3.2cm] (bytes) {载荷事实\\长度、终点、校验};
\node[stage,minimum width=3.2cm,right=of bytes] (layout) {布局事实\\入口、零区、栈};
\node[stage,minimum width=3.2cm,right=of layout] (state) {现场事实\\寄存器、SP、返回地址};
\node[stage,minimum width=3.3cm,below=14mm of layout] (commit) {唯一提交动作\\PC 写入任务入口};
\draw[flow] (bytes) -- (layout);
\draw[flow] (layout) -- (state);
\draw[flow] (bytes.south) |- (commit.west);
\draw[flow] (layout) -- (commit);
\draw[flow] (state.south) |- (commit.east);
\end{tikzpicture}
\end{center}
\diagramnote{三组事实都成立，才允许执行控制权提交。PC 写入不是又一次“检查”，而是使任务
开始取指的不可混淆边界。}

\subsection{用装入记录快照检查一次正常传输}

下面选择四个时刻观察装入记录。设载荷长度为 \code{0x2100}，起点为
\code{0x00100000}。

\begin{longtable}{|L{0.19\linewidth}|L{0.21\linewidth}|L{0.22\linewidth}|L{0.26\linewidth}|}
\hline
观察时刻 & 下一写入地址 & 剩余字节 & 已经成立的事实 \\ \hline
头部通过后 & \code{0x00100000} & \code{0x2100} & 范围与入口合法，尚无载荷字节 \\ \hline
第 1 字节后 & \code{0x00100001} & \code{0x20ff} & 起始字节已写入并纳入校验 \\ \hline
第 64 字节后 & \code{0x00100040} & \code{0x20c0} & 第一块可向发送端确认 \\ \hline
最后字节后 & \code{0x00102100} & \code{0} & 长度满足，等待最终校验 \\ \hline
\end{longtable}

两个数值始终满足不变量：
\[
\texttt{destination\_next}-\texttt{load\_address}
=\texttt{image\_bytes}-\texttt{payload\_remaining}.
\]
左边是已经写入的地址跨度，右边是已经消费的字节数。若更新时漏掉递增或递减，其中一边会
立刻失配。固件可在调试版本中每轮检查该式。

\subsection{第 37 个字节损坏时哪些字节可见}

假设传输没有丢失长度，但第 37 个字节从 \code{0x11} 变成 \code{0x10}。装入器会继续收满
\code{0x2100} 字节，最终校验不匹配，状态转为 \code{REJECTED}。此时 RAM 中确实有一份
长度完整但内容错误的映像。

固件不必在发现校验失败时逐字节回滚，因为任务 PC 尚未提交。它必须保证两点：不把记录
转为 \code{READY}；下一次装入重新覆盖可能使用的范围。若机器要在多个已装入映像之间
保留旧版本，直接覆盖就不再足够，需要独立暂存区或双缓冲；本章只有一项任务，因此不引入
该结构。

\subsection{第 37 个字节之后串行溢出}

若控制器置 \code{RX\_OVERRUN}，固件知道至少一个字节丢失，却不知道具体有几个。此时
\code{destination\_next} 只表示已经消费的字节数，不能对应发送端位置。固件立即进入
\code{REJECTED}，丢弃直到下一次明确的传输同步标记，再请求整包重发。

仅把 \code{payload\_remaining} 继续减到零会造成错误对齐：后续字节向前填补缺口，头部
长度仍可能满足，但内容位置全部错位。校验大概率会失败，却要等整包结束才能发现。尽早
利用硬件溢出状态可以缩短失败路径。

\subsection{零填充区为什么不随载荷发送}

任务常需要一片初始值全为零的可写区域，例如计数器数组或暂存缓冲。如果发送端把几千个零
也逐字节发送，既增加传输时间，又没有增加信息。头部用 \code{zero\_bytes} 表达“映像末尾
之后这段范围必须被固件写零”。

这一做法建立了两类来源不同的任务字节：

\begin{longtable}{|L{0.24\linewidth}|L{0.26\linewidth}|L{0.36\linewidth}|}
\hline
区域 & 初值来源 & 校验与建立规则 \\ \hline
映像区 & 外部载荷 & 每个字节参与传输校验 \\ \hline
零填充区 & 固件写零 & 范围由已验证头部决定，不占载荷 \\ \hline
栈区 & 固件初始化 & 预置返回地址，其余按启动策略清零 \\ \hline
\end{longtable}

零填充必须发生在校验通过之后或之前均可，但在状态进入 \code{READY} 前必须完成。若清零
中途发生复位，启动流程重新开始，不能沿用旧的 \code{READY} 标记。

\subsection{为什么不能把整个 RAM 都清零}

全清零看似简单，却会覆盖固件工作区，包括当前装入记录和固件栈。即使先避开工作区，清零
整片 RAM 也会让启动时间随容量增长。按头部和固定布局只初始化任务真正依赖的区域，使
写入范围与所有权一致。

\subsection{任务栈如何支持一次真实函数调用}

求和核心本身是叶函数，但正式任务通常会调用子程序。设任务在写结果前调用位于
\code{0x00100100} 的格式化例程，调用指令地址为 \code{0x00100018}，顺序后继为
\code{0x0010001c}。调用前：

\begin{lstlisting}
PC = 0x00100018
SP = 0x001efffc
MEM32[0x001efffc] = 0x00000300   // outer return
\end{lstlisting}

\code{CALL} 先把 SP 降到 \code{0x001efff8}，写入内部返回地址
\code{0x0010001c}，再进入子程序：

\begin{lstlisting}
PC = 0x00100100
SP = 0x001efff8
MEM32[0x001efff8] = 0x0010001c   // return to task
MEM32[0x001efffc] = 0x00000300   // return to firmware
\end{lstlisting}

子程序的 \code{RET} 先消费 \code{0x0010001c}，恢复到任务内部，SP 回到
\code{0x001efffc}。任务最外层 \code{RET} 再消费 \code{0x00000300}，返回固件。两个
返回地址按后进先出次序工作。

\begin{center}
\begin{tikzpicture}[x=1cm,y=0.7cm]
\draw[BookRule,thick] (0,0) rectangle (8,4.4);
\draw[BookRule] (0,1.1) -- (8,1.1);
\draw[BookRule] (0,2.2) -- (8,2.2);
\draw[BookRule] (0,3.3) -- (8,3.3);
\node[anchor=west] at (0.2,3.85) {\code{0x001f0000}：栈区上界};
\node[anchor=west] at (0.2,2.75) {\code{0x001efffc}：返回固件 \code{0x00000300}};
\node[anchor=west] at (0.2,1.65) {\code{0x001efff8}：返回任务 \code{0x0010001c}};
\node[anchor=west] at (0.2,0.55) {更深调用与局部保存继续向低地址扩展};
\draw[flow] (9.3,1.65) -- (8.05,1.65);
\node[anchor=west] at (9.4,1.65) {子程序中的 SP};
\end{tikzpicture}
\end{center}
\diagramnote{返回地址的顺序来自 SP 的移动和 RAM 中的布局，不来自处理器对“函数层次”的
理解。破坏任一栈字都会改变后续控制流。}

\subsection{谁负责保证栈不越界}

当前处理器只检查地址是否位于整个 RAM，不知道任务栈下界
\code{0x001ec000}。若调用过深，SP 可以越过下界并覆盖其他任务数据，只要地址仍在 RAM，
硬件就会接受。固件给出 16 KiB 空间并不能强制任务停在边界内。

可以让编写任务的人证明最大栈深，或在每次函数入口加入软件检查；两者仍依赖任务合作。
真正的隔离需要地址访问权限，后续才会从任务相互覆盖的失败中推导。

\subsection{移交失败时由谁收束状态}

在 PC 提交之前，固件是唯一执行者，也是装入记录、任务目标区和串行协议的所有者。任何
失败都由固件把记录转为 \code{REJECTED}，发送错误码，并回到等待新头部状态。

PC 提交之后，固件不再执行。任务拥有当前寄存器与任务栈，直到 \code{RET}。此后出现的
无限循环、任意 RAM 写入和栈破坏，不能再由“装入器检查一次”解决。阶段不同，能采取行动
的主体也不同。

\begin{longtable}{|L{0.20\linewidth}|L{0.24\linewidth}|L{0.28\linewidth}|L{0.18\linewidth}|}
\hline
阶段 & 当前执行者 & 可修改的关键状态 & 失败收束 \\ \hline
接收头部 & 固件 & 装入记录、串行槽 & 拒绝传输 \\ \hline
接收载荷 & 固件 & 任务 RAM、校验状态 & 拒绝并重传 \\ \hline
建立现场 & 固件 & 任务栈、参数寄存器 & 不提交 PC \\ \hline
任务运行 & 任务 & 寄存器及全部可写物理 RAM & 合作返回或整机错误入口 \\ \hline
固定返回 & 固件 & 固件栈、发送状态 & 报告结果并等待 \\ \hline
\end{longtable}

这张所有权表给出了本章最重要的时间边界：同一片 RAM 在不同阶段由不同代码解释和修改，
但硬件尚未实施所有权。当前所有权是分析软件责任的工具，不是安全保证。

\section{把固件主循环连成完整控制程序}

前面已经逐个定义装入记录、串行寄存器和任务现场。现在可以给出不隐藏阶段的固件主循环。
伪代码使用函数名表达已经解释过的局部动作，不涉及构建工具或源文件位置。

\begin{lstlisting}
reset_entry:
  initialize firmware SP
  clear serial protocol state
  reset load_record to WAITING

wait_for_transfer:
  header = receive_exactly(32)
  if serial overrun: report SERIAL_OVERRUN; goto wait_for_transfer
  if not validate_header(header):
    report HEADER_INVALID
    goto wait_for_transfer

  initialize load_record from header
  while payload_remaining != 0:
    b = receive_one_byte()
    if serial overrun:
      reject SERIAL_OVERRUN
      goto discard_and_resync
    store b at destination_next
    update checksum and counters
    acknowledge each completed 64-byte block

  load_record.state = VERIFYING
  if checksum_so_far != expected_checksum:
    report CHECKSUM_MISMATCH
    goto wait_for_transfer

  zero requested region
  build task stack and arguments
  assert handoff invariants
  load_record.state = READY
  handoff_to_task()        // PC commit occurs here
\end{lstlisting}

\code{handoff\_to\_task} 正常情况下不返回到下一行。它写入任务寄存器和 SP，最后通过受控
跳转提交 PC。任务执行最外层 \code{RET} 时进入另一个固定入口
\code{firmware\_return}，而不是回到调用点。

\subsection{为什么返回入口不能沿用装入函数的普通栈帧}

装入函数的局部变量位于固件栈。移交任务时，SP 已切到任务栈；任务也可以改写所有通用
寄存器。若返回入口试图像普通函数那样从旧 SP 取局部变量，它会读错位置。固定入口必须先
把 SP 设为固件栈顶，再通过全局固定地址找到装入记录。

\begin{lstlisting}
firmware_return:
  SP = 0x00200000
  load_record.state = FINISHED
  result = read32(0x00102020)
  send_result_frame(result)
  reset load_record to WAITING
  goto wait_for_transfer
\end{lstlisting}

这里 \code{0x00200000} 是 RAM 上界，固件实际第一项栈内容写在其下方。固件工作区位于
\code{0x001f0000--0x001fffff}，任务栈止于 \code{0x001effff}，两者按软件布局分离。

\subsection{结果帧为什么还要带状态}

只发送四个结果字节，外部端无法区分正常结果、错误码和上一次传输残留。固件发送一个短帧：

\begin{lstlisting}
result_frame:
  magic      = "RES1"
  status     = OK or an error code
  value      = valid only when status == OK
  checksum   = checksum of preceding fields
\end{lstlisting}

本次成功帧中的 value 为 32 位整数 21，在线上按
\code{15 00 00 00} 发送。外部端只有在完整接收并验证结果帧后，才把任务标记为成功。
固件写 \code{TXDATA} 只是提交一个待发送字节，不等于外部端已经收到整帧。

\subsection{从复位到等待下一任务的状态图}

\begin{center}
\begin{tikzpicture}[node distance=9mm and 13mm]
\node[stage,minimum width=2.8cm] (reset) {复位初始化};
\node[stage,minimum width=2.8cm,right=of reset] (wait) {等待头部};
\node[stage,minimum width=2.8cm,right=of wait] (recv) {接收载荷};
\node[stage,minimum width=2.8cm,below=13mm of recv] (verify) {验证与建栈};
\node[stage,minimum width=2.8cm,left=of verify] (run) {任务运行};
\node[stage,minimum width=2.8cm,left=of run] (report) {报告结果};
\node[stage,minimum width=2.8cm,below=13mm of wait,fill=BookWarm] (reject) {报告拒绝};
\draw[flow] (reset)--(wait);
\draw[flow] (wait)--node[above,font=\scriptsize]{合法头部}(recv);
\draw[flow] (recv)--node[right,font=\scriptsize]{载荷到齐}(verify);
\draw[flow] (verify)--node[above,font=\scriptsize]{提交 PC}(run);
\draw[flow] (run)--node[above,font=\scriptsize]{合作返回}(report);
\draw[flow] (report.west)-| ++(-8mm,13mm) |- (wait.west);
\draw[->,>=Latex,thick,draw=BookMuted] (wait)--(reject);
\draw[->,>=Latex,thick,draw=BookMuted] (recv.west)-|(reject.east);
\draw[->,>=Latex,thick,draw=BookMuted] (verify.north)-|(reject.east);
\draw[flow] (reject)--(wait);
\end{tikzpicture}
\end{center}
\diagramnote{固件能控制的转换都发生在“任务运行”之外。进入任务后，唯一正常出边是任务主动
返回；当前机器没有从运行状态强制到达报告或等待状态的事件。}

\subsection{状态图中的终止条件}

等待头部没有时间上限；机器可以长期停在那里。接收载荷若外部端中途停止，当前协议也会
无限等待，因为机器尚无可编程计时器。若希望超时拒绝，必须增加能产生时间事件的硬件和
相应软件状态。本章不借用尚未加入的能力。

任务运行同样可能没有终点。状态图不是承诺所有路径最终回到等待，而是准确标出哪些转换
需要外部合作。正式模型必须允许停留在 \code{RECEIVING} 或任务运行状态。

\subsection{一次完整成功运行中的关键数值}

\begin{longtable}{|L{0.23\linewidth}|L{0.28\linewidth}|L{0.35\linewidth}|}
\hline
\rowcolor{BookTableHead}
时刻 & 关键状态 & 可据此推出的事实 \\ \hline
复位释放 & \code{PC=0x00000000} & 下一笔取指来自启动存储 \\ \hline
头部通过 & \code{state=RECEIVING} & 地址与长度安全，内容尚未验证 \\ \hline
载荷到齐 & \code{remaining=0} & 已写 \code{0x2100} 字节 \\ \hline
校验通过 & \code{state=VERIFYING} & 传输内容与发送端一致 \\ \hline
现场完成 & \code{SP=0x001efffc} & 栈顶含固件返回地址 \\ \hline
移交提交 & \code{PC=0x00100000} & 任务开始取指 \\ \hline
第四轮后 & \code{r4=21,r1=0} & 循环将退出 \\ \hline
结果写回 & \code{MEM32[0x00102020]=21} & 处理器可读取结果 \\ \hline
任务返回 & \code{PC=0x00000300} & 固件重新获得控制 \\ \hline
结果帧到达 & 外部校验通过 & 本次任务对外完成 \\ \hline
\end{longtable}

这十个时刻覆盖了硬件起点、协议接受、内容验证、执行移交、计算效果、控制返回和外部可见。
任意相邻两行之间都存在不能省略的动作，因此不能把它们压缩成“开机后运行程序并输出结果”。

\subsection{本章方案的成本}

最小方案追求边界清楚，不追求吞吐最大。处理器轮询串行接口时不能做其他工作；每个任务都
要完整传输；任务运行时固件完全停止；结果只存在 RAM；错误通常通过复位收束。这些成本
不是待补的一串现代名词，而是当前路径中可以指认的位置。

\begin{longtable}{|L{0.23\linewidth}|L{0.29\linewidth}|L{0.34\linewidth}|}
\hline
\rowcolor{BookTableHead}
当前成本 & 在路径中的原因 & 只有何时才值得改进 \\ \hline
装入时处理器忙等 & 单槽串行接口只提供轮询状态 & 输入速度或并行工作成为问题 \\ \hline
启动需完整传输 & RAM 断电丢失且任务不在 ROM & 需要持久保存多个任务 \\ \hline
不能同时运行别的任务 & 只有一份当前处理器现场 & 第二任务需要取得处理器 \\ \hline
无限循环无法收束 & 没有异步时间事件 & 必须限制任务运行区间 \\ \hline
任意 RAM 写可破坏固件 & 物理地址没有权限 & 不可信任务需要隔离 \\ \hline
\end{longtable}

下一章首先处理第三行：在不要求任务重新从入口开始的条件下，让第二项任务获得处理器。其余
问题继续保留，避免一次引入尚未由失败要求的结构。

\section{沿着处理器状态走完求和任务}

任务已经获得处理器。现在用具体地址和数值逐条检查执行结果。任务开始时输入内存为：

\begin{lstlisting}
MEM32[0x00102000] = 7
MEM32[0x00102004] = -2
MEM32[0x00102008] = 11
MEM32[0x0010200c] = 5
MEM32[0x00102020] = 0
\end{lstlisting}

\subsection{初始化指令}

第一笔任务取指从 \code{0x00100000} 到达 RAM。互连不关心这是任务代码，只因地址匹配而
选择 RAM。返回字节译码为 \code{MOVI r4,0}。指令提交后：

\begin{lstlisting}
PC: 0x00100000 -> 0x00100004
r4: 0x00000000 -> 0x00000000
other architectural state: unchanged
\end{lstlisting}

虽然 \code{r4} 的数值碰巧没有变化，写入动作仍由指令语义决定。若固件没有按约定清零
\code{r4}，这条指令也会建立相同起点。

\subsection{第一轮：从字节到部分和}

在 \code{0x00100004}，\code{LOAD r5,[r0+0]} 读取
\code{0x00102000--03}。事务成功后 \code{r5=7}，PC 进入
\code{0x00100008}。\code{ADD r4,r5} 读取旧值 \(0\) 和 \(7\)，产生候选结果 \(7\)，
提交后 \code{r4=7}。

接下来的两条 \code{ADDI} 分别改变指针和计数：

\begin{lstlisting}
r0: 0x00102000 -> 0x00102004
r1: 4          -> 3
\end{lstlisting}

\code{BNEZ r1,0x00100004} 读取 \code{r1=3}，条件成立，因此下一 PC 不是顺序地址
\code{0x00100018}，而是循环入口 \code{0x00100004}。

\subsection{四轮状态表}

每轮在条件分支提交后，关键状态如下：

\begin{longtable}{|L{0.12\linewidth}|L{0.22\linewidth}|L{0.16\linewidth}|L{0.16\linewidth}|L{0.18\linewidth}|}
\hline
\rowcolor{BookTableHead}
轮次 & 本轮读地址 & 读出值 & 部分和 \code{r4} & 分支后的 \code{r1} \\ \hline
1 & \code{0x00102000} & \(7\) & \(7\) & \(3\) \\ \hline
2 & \code{0x00102004} & \(-2\) & \(5\) & \(2\) \\ \hline
3 & \code{0x00102008} & \(11\) & \(16\) & \(1\) \\ \hline
4 & \code{0x0010200c} & \(5\) & \(21\) & \(0\) \\ \hline
\end{longtable}

第四轮后，\code{BNEZ} 不采用分支目标，PC 顺序进入 \code{0x00100018}。此时
\code{r0=0x00102010}，正好指向数组末尾之后的位置。这个地址没有被解引用；它只表示所有
元素已经消费。若循环多执行一轮，处理器会读取数组外字节，而当前机器没有边界检查替任务
阻止这一错误。

\subsection{结果写入的可见边界}

\code{STORE r4,[r2+0]} 形成地址 \code{0x00102020}，写数据为
\code{0x00000015}，宽度为 4。RAM 在事务完成时把四个字节更新为
\code{15 00 00 00}。在写响应返回前，任务不能把该存储指令视为完成。

写入完成表示后续处理器读取能够取得新值。它不表示结果跨断电保存，因为目标是 RAM；也不
表示外部发送端已经看到结果，因为串行发送尚未发生。可见、传达和持久是三个不同边界。

\begin{center}
\begin{tikzpicture}[node distance=10mm and 8mm]
\node[stage,minimum width=2.5cm] (reg) {寄存器结果\\\code{r4=21}};
\node[stage,minimum width=2.8cm,right=of reg] (ram) {RAM 写完成\\处理器可读};
\node[stage,minimum width=2.7cm,right=of ram] (tx) {串行发送\\外部端收到};
\node[stage,minimum width=2.8cm,right=of tx,fill=BookWarm] (power) {断电后保留\\本章不保证};
\draw[flow] (reg) -- node[above,font=\scriptsize]{STORE} (ram);
\draw[flow] (ram) -- node[above,font=\scriptsize]{固件报告} (tx);
\draw[->,>=Latex,dashed,thick,draw=BookMuted] (tx) -- (power);
\end{tikzpicture}
\end{center}
\diagramnote{前三个方框是不同层次的事实。最右侧使用虚线，因为 RAM 和串行发送都不能提供
断电持久性。}

\subsection{任务怎样返回固件}

执行到 \code{0x0010001c} 时，任务发出 \code{RET}。处理器必须先从
\code{SP=0x001efffc} 读取 4 字节返回地址。RAM 返回 \code{0x00000300} 后，处理器把
SP 增加到 \code{0x001f0000}，并把 PC 改为 \code{0x00000300}。这两个状态更新共同完成
返回。

\subsection{返回前后快照}

\begin{longtable}{|L{0.26\linewidth}|L{0.27\linewidth}|L{0.27\linewidth}|}
\hline
\rowcolor{BookTableHead}
状态 & \code{RET} 开始前 & \code{RET} 完成后 \\ \hline
\code{PC} & \code{0x0010001c} & \code{0x00000300} \\ \hline
\code{SP} & \code{0x001efffc} & \code{0x001f0000} \\ \hline
\code{MEM32[0x001efffc]} & \code{0x00000300} & 内容仍在，但已不属于有效栈 \\ \hline
\code{MEM32[0x00102020]} & \(21\) & \(21\) \\ \hline
\end{longtable}

SP 上升不会自动擦除旧返回地址。栈的有效范围由 SP 与约定共同决定，旧字节可以继续留在
RAM 中。后续压栈可能覆盖它。

\subsection{固件返回入口能依赖什么}

固件不能假定任务保留了自己的临时寄存器，因为本章没有要求任务这样做。固定返回入口只
依赖两项事实：PC 已指向 \code{0x00000300}，结果位于约定的 RAM 地址。它重新建立自己的
SP，读取结果，等待串行发送端可用，然后按四个字节发送结果或发送文本表示。

\begin{lstlisting}
firmware_return:
  SP = FIRMWARE_STACK_TOP
  result = MEM32[0x00102020]
  for each byte in encode(result):
    repeat until STATUS has TX_READY
    write8(TXDATA, byte)
  enter firmware_wait_loop
\end{lstlisting}

固件必须先重建 SP，再进行函数调用或压栈。任务返回时的 SP 已在任务栈上界，不能当作固件
栈继续使用。两片栈位于同一 RAM，却属于不同执行阶段的软件约定。

\section{正常路径的端到端时间线}

前面分别建立了器件和局部事务，现在可以把它们综合成一次完整运行。下图只使用已经解释过
的对象：

\begin{center}
\begin{tikzpicture}[x=2.25cm,y=0.78cm]
\node[font=\footnotesize] at (0,7.5) {外部端};
\node[font=\footnotesize] at (1,7.5) {串行控制器};
\node[font=\footnotesize] at (2,7.5) {固件};
\node[font=\footnotesize] at (3,7.5) {RAM};
\node[font=\footnotesize] at (4,7.5) {任务};
\foreach \x in {0,1,2,3,4} {\draw[BookRule] (\x,7.2) -- (\x,0.1);}
\draw[flow] (0,6.8) -- node[above,font=\scriptsize]{头部字节} (1,6.8);
\draw[flow] (2,6.35) -- node[above,font=\scriptsize]{轮询读取} (1,6.35);
\draw[flow] (2,5.8) -- node[above,font=\scriptsize]{验证头部} (3,5.8);
\draw[flow] (0,5.2) -- node[above,font=\scriptsize]{载荷字节} (1,5.2);
\draw[flow] (2,4.7) -- node[above,font=\scriptsize]{逐字节写入} (3,4.7);
\draw[flow] (2,4.1) -- node[above,font=\scriptsize]{校验、清零、建栈} (3,4.1);
\draw[flow] (2,3.5) -- node[above,font=\scriptsize]{提交入口 PC} (4,3.5);
\draw[flow] (4,2.9) -- node[above,font=\scriptsize]{读取输入} (3,2.9);
\draw[flow] (4,2.3) -- node[above,font=\scriptsize]{写结果 21} (3,2.3);
\draw[flow] (4,1.7) -- node[above,font=\scriptsize]{RET 到固定入口} (2,1.7);
\draw[flow] (2,1.1) -- node[above,font=\scriptsize]{写 TXDATA} (1,1.1);
\draw[flow] (1,0.55) -- node[above,font=\scriptsize]{结果字节} (0,0.55);
\end{tikzpicture}
\end{center}
\diagramnote{时间从上向下。载荷进入 RAM、任务获得 PC、结果进入 RAM、结果到达外部端分别是
四个不同事件；前一事件不会自动包含后一事件。}

\subsection{把每个边界写成一句可检验的事实}

\begin{enumerate}
\item 头部接受：字段合法且后续计算不会越过任务区域；
\item 载荷接收完成：规定数量的字节已经写入 RAM；
\item 映像验证完成：校验通过、零区和栈均已建立；
\item 控制权移交：PC 已提交任务入口，下一次取指来自任务；
\item 计算完成：结果存储事务已得到 \code{OK}；
\item 任务返回：PC 和 SP 已按 \code{RET} 规则更新；
\item 外部可见：串行发送完成，外部端实际收到结果字节。
\end{enumerate}

如果只说“任务运行成功”，就无法定位失败究竟发生在传输、装入、移交、计算、返回还是
报告阶段。正式系统中的请求、完成和持久化边界更复杂，但分析方法从这台小机器开始已经
相同：先命名对象状态，再指出哪个事件改变了它。

\section{错误路径揭示了当前机器缺少什么}

正常示例证明最小方案可行，错误路径则说明它的边界。下面每种错误都对应一个具体检测者，
不能用“系统会处理”概括。

\subsection{头部损坏：固件可以在移交前拒绝}

若 \code{magic} 错误、长度越界或入口未对齐，固件仍掌握 PC，因此可以停止接收、通过
串行发送错误码并回到等待头部状态。此时 RAM 中即使残留部分字节，也没有被执行。

错误处理的关键不是把 RAM 恢复成全零，而是不提交任务入口。下一次合法装入会覆盖规定
范围；若固件需要避免旧数据泄漏，可以在重新使用前清零，但这属于额外策略。

\subsection{串行溢出：控制器只能报告已经丢失}

若发送端速度超过固件轮询速度，新字节会覆盖单槽接收寄存器，控制器随即置溢出位
\code{RX\_OVERRUN}。固件无法由该位推断丢失字节的内容和位置，只能拒绝整次传输并
要求外部端重发。校验和可以检测内容不一致，却不能凭空恢复缺失字节。

增加接收 FIFO、流量控制或 DMA 都能改变吞吐边界，但当前问题只需要一个最小可靠方案，
所以采用“检测溢出并重传”的规则。后续若设备速率成为主要矛盾，再引入更复杂硬件。

\subsection{未映射地址：互连报告错误，任务没有恢复机制}

若任务错误地读取 \code{0x20000000}，互连返回 \code{UNMAPPED}。教学处理器把失败地址、
访问方向和当时 PC 写入固定固件错误槽，然后跳到固件错误入口。这使机器至少不会把任意
电平当作成功数据。

本章尚未为不同任务建立独立错误现场，也没有把错误转换成可返回给应用的信号。故障后的
最小处理是报告并复位。错误可检测不等于错误可恢复。

\subsection{任务写坏固件工作区：软件约定无法强制隔离}

任务执行：

\begin{lstlisting}
STORE r4, [0x001f0000]
\end{lstlisting}

时，地址合法地落在 RAM 范围。互连和 RAM 都会完成写入，因为它们不知道该范围按约定属于
固件。任务随后返回时，固件的装入记录或栈可能已经损坏。这不是地址互连故障，而是当前
机器缺少访问权限边界。

\subsection{栈指针损坏：返回地址的存在也不足以保证返回}

若任务把 SP 改成 \code{0x00102000} 后执行 \code{RET}，处理器会把第一个输入整数
\code{0x00000007} 当作返回 PC。该地址未按指令对齐或落在无效区域，机器随后进入错误
入口。固件预置返回地址只能保证合作任务有一条返回路径，不能保证任意代码使用它。

\subsection{无限循环：没有外部事件就无法取回处理器}

任务可能执行：

\begin{lstlisting}
repeat:
  ADDI r4, 1
  BNEZ r4, repeat
\end{lstlisting}

只要指令和 RAM 访问合法，处理器就会继续执行任务。串行控制器即使收到新命令，当前代码
也没有读取它。时钟边沿只能推动指令前进，并不会把 PC 自动改回固件。机器此时既没有可编程
时钟事件，也没有异步入口，固件无法重新获得控制。

\begin{problemframe}
\textbf{这一失败是后续演进的关键证据。} “任务最终会执行 \code{RET}”是合作约定，不是
机器保证。若希望运行多个任务或限制单个任务占用处理器的时间，必须先回答怎样保存正在
执行的现场，再回答怎样在任务不合作时进入受控代码。
\end{problemframe}

\subsection{这台机器的软件模型究竟是什么}

到本章末，程序员面对的不是一堆孤立器件，而是一组边界清楚的契约。

\subsection{地址契约}

一个 32 位物理地址经过互连选择启动存储、RAM、串行控制器或错误响应。读写宽度和方向是
事务的一部分。地址只标识位置，不表达对象身份、权限或所有权。

\subsection{执行契约}

PC 指出下一条指令，寄存器与 RAM 保存继续计算所需的状态。普通指令顺序推进，分支、调用
和返回按明确规则改变 PC。只有已经提交的状态才属于下一条指令的输入。

\subsection{启动契约}

复位先进入固件。固件验证任务映像，建立内存布局、参数和栈，最后提交任务入口 PC。任务
执行 \code{RET} 才会沿预置栈内容回到固定固件入口。

\subsection{设备契约}

串行控制器把异步到达的线信号变成带状态位的字节槽。固件通过轮询消费字节，并负责在溢出
或校验失败时拒绝整次传输。控制器不理解任务格式。

\section{为什么这仍然不是操作系统}

固件已经做了不少软件工作：它初始化器件、接收字节、检查范围、装入任务、建立初始现场并
报告结果。但“存在启动软件”不等于“已经有操作系统”。当前方案只服务于一条预先约定的
运行路径，没有建立长期存在的任务身份和资源管理规则。

\begin{longtable}{|L{0.24\linewidth}|L{0.29\linewidth}|L{0.35\linewidth}|}
\hline
\rowcolor{BookTableHead}
当前已经具备 & 依靠什么实现 & 当前仍未具备 \\ \hline
确定的复位入口 & 复位状态与启动存储 & 多个任务的身份和生命周期 \\ \hline
任务字节装入 RAM & 串行协议与固件装入记录 & 文件、目录和持久命名 \\ \hline
一次明确的控制权移交 & 初始寄存器、栈和 PC 提交 & 暂停后恢复任意任务现场 \\ \hline
合作式返回 & 预置返回地址与 \code{RET} & 强制收回处理器 \\ \hline
物理地址访问 & 地址互连 & 地址隔离和访问权限 \\ \hline
错误检测与复位 & 互连错误和固件入口 & 面向任务的独立错误恢复 \\ \hline
\end{longtable}

表的左列是本章已经逐步建立并走读过的能力，右列则是具体失败留下的空缺。后续章节不会
直接罗列现代操作系统模块，而会从这些空缺中一次选择一个进行演进。

\subsection{本章复盘：从一个要求到一台可运行机器}

最初的要求只有“计算四个整数之和”。为了让这句话在真实机器上成立，我们依次得到以下
因果链：

\begin{enumerate}
\item 计算必须跨动作保留事实，因此需要寄存器、PC、SP 和可写状态；
\item 第一条指令必须跨断电存在，因此需要启动存储；
\item 运行数据必须可变，因此需要 RAM；
\item 多个器件不能同时响应同一请求，因此需要地址互连和明确事务；
\item 上电状态不可靠，因此需要时钟边界、复位状态和固定入口；
\item 外部任务不能凭空进入 RAM，因此需要串行控制器和传输协议；
\item 任意字节不能直接执行，因此需要范围、校验和入口验证；
\item 任务不能继承固件偶然状态，因此需要初始寄存器、独立栈和移交顺序；
\item 任务完成后必须有控制流去向，因此需要预置返回地址和固定固件入口。
\end{enumerate}

这条链不是整机启动的口号，而是每一步都有对象、字段和状态变化的运行记录。求和任务的
最终结果为 \(21\)，RAM 中的结果字节为 \code{15 00 00 00}，外部端只有在串行发送完成后
才能看到它。

\section{为什么四轮结束时一定得到 21}

逐条 trace 能证明本次输入的实际运行结果，还可以从循环状态中提炼一个每轮都成立的关系。
设原数组为 \(a_0,a_1,a_2,a_3\)，总元素数为 \(n=4\)，已经处理的元素数为 \(k\)。每次
进入 \code{LOAD} 前，程序应满足：
\[
\begin{aligned}
\texttt{r0} &= \texttt{array\_base}+4k,\\
\texttt{r1} &= n-k,\\
\texttt{r4} &= \sum_{i=0}^{k-1}a_i.
\end{aligned}
\]
这个关系称为循环不变量，因为初始化建立它，每一轮保持它，退出条件再利用它推出结果。

\subsection{初始化为什么建立不变量}

第一轮之前 \(k=0\)。固件令 \code{r0=array\_base}、\code{r1=4}，第一条任务指令令
\code{r4=0}。空前缀的和定义为 0，因此三个等式全部成立。这里能看出任务代码和启动约定
共同建立初态：若固件传错 \code{r0} 或 \code{r1}，单靠 \code{MOVI r4,0} 不能修复。

\subsection{一轮怎样保持不变量}

假设进入第 \(k\) 轮时关系成立。\code{LOAD} 从
\(\texttt{array\_base}+4k\) 读取 \(a_k\)，\code{ADD} 把它加入 \code{r4}；
\code{ADDI r0,4} 令指针对应 \(k+1\)，\code{ADDI r1,-1} 令剩余数变成
\(n-(k+1)\)。于是下一次进入循环时：
\[
\texttt{r4}=\sum_{i=0}^{k}a_i,
\]
正好是把 \(k\) 替换为 \(k+1\) 后的不变量。

\subsection{退出条件怎样推出最终结果}

分支不采用循环目标时 \code{r1=0}。由 \(\texttt{r1}=n-k\) 得 \(k=n=4\)，因此
\[
\texttt{r4}=a_0+a_1+a_2+a_3=7-2+11+5=21.
\]
\code{STORE} 把该值写到结果槽。这个推导还揭示零长度问题：当前代码在第一次分支检查前
就执行装入，只有 \(n>0\) 时初始化到第一次循环入口的路径才满足预期使用条件。

\begin{center}
\begin{tikzpicture}[node distance=11mm and 16mm]
\node[stage,minimum width=3.3cm] (init) {初始化\\\(k=0,\ r4=0\)};
\node[stage,minimum width=3.6cm,right=of init] (body) {处理 \(a_k\)\\指针前移、计数减少};
\node[stage,minimum width=3.3cm,right=of body] (test) {检查\\\(n-(k+1)\)};
\node[stage,minimum width=3.4cm,below=14mm of body] (done) {退出\\\(k=n,\ r4=\sum a_i\)};
\draw[flow] (init)--(body);
\draw[flow] (body)--(test);
\draw[flow] (test.north) -- ++(0,8mm) -| node[above,font=\scriptsize,pos=0.55]{仍有元素} (body.north);
\draw[flow] (test.south) |- node[right,font=\scriptsize,pos=0.3]{计数为零} (done.east);
\end{tikzpicture}
\end{center}
\diagramnote{机器 trace 给出每步具体值，不变量说明这些值为何对任意满足启动契约的非空数组
都沿同一关系推进。}

\subsection{装入一项任务需要多少次数据移动}

载荷长度 \code{0x2100} 等于 8448 字节，也恰好是 132 个 64 字节块。对每个字节，串行
控制器先把外部位收进 \code{RXDATA}；处理器读取一次 \code{RXDATA}；随后处理器向 RAM
写入一次。仅计算有效数据移动，就有 8448 笔设备读和 8448 笔 RAM 写。

状态轮询次数不固定。如果每个字节到达时固件恰好检查一次，需要至少 8448 笔
\code{STATUS} 读；若固件检查十次才等到一个字节，状态读会增加到约 84480 笔。轮询的
成本来自“没有数据时仍要主动读取状态”。

\begin{longtable}{|L{0.23\linewidth}|L{0.24\linewidth}|L{0.37\linewidth}|}
\hline
动作 & 本次最少次数 & 次数由什么决定 \\ \hline
读取 \code{RXDATA} & 8448 & 载荷字节数 \\ \hline
写入任务 RAM & 8448 & 载荷字节数 \\ \hline
读取 \code{STATUS} & 不少于 8448 & 到达间隔与轮询频率 \\ \hline
发送块确认 & 132 次确认帧 & 64 字节分块规则 \\ \hline
零填充写入 & 256 字节 & \code{zero\_bytes=0x100} \\ \hline
建立返回地址 & 1 笔 4 字节写 & 启动栈约定 \\ \hline
\end{longtable}

这个计数没有把取指事务算入，因为固件执行每条轮询和搬运指令还要取指。即使不建立精确
周期模型，也能看出装入成本至少与载荷长度线性增长。DMA 可以减少逐字节处理器搬运，但要
引入描述符、缓冲所有权和完成事件；当前规模没有要求它。

\subsection{确认帧不是免费的}

每 64 字节等待一次确认会限制连续突发长度，也增加往返时间。若确认帧占 8 个串行字节，
132 次确认会额外发送 1056 字节。把块增大到 256 字节可减少确认次数，却会让出错后重传
更多数据，并要求控制器或协议容纳更长未确认区间。

因此块大小不是任意常数。它在状态开销、错误恢复粒度和缓冲需求之间取舍。本章固定为
64 字节，只为使单槽轮询方案具有明确背压点。

\subsection{三个快照足以区分“装入”“运行”和“返回”}

\subsection{快照一：映像已经在 RAM，但固件仍运行}

\begin{lstlisting}
load_record.state = VERIFYING
PC = 0x00000280
SP = 0x001ffec0
MEM[0x00100000..0x001020ff] = received image
task registers = not established
\end{lstlisting}

这时读取任务区域能看到完整字节，却不能说任务已运行。PC 仍在启动存储，寄存器和栈也仍
属于固件阶段。

\subsection{快照二：第一条任务指令尚未提交}

\begin{lstlisting}
load_record.state = RUNNING
PC = 0x00100000
SP = 0x001efffc
r0 = 0x00102000
r1 = 4
r2 = 0x00102020
\end{lstlisting}

PC 已指向任务入口，下一笔取指将访问 RAM。第一条 \code{MOVI} 尚未提交，因此这个快照是
控制权移交后的初始任务现场。

\subsection{快照三：任务返回入口刚获得控制}

\begin{lstlisting}
PC = 0x00000300
SP = 0x001f0000
r4 = 21
MEM32[0x00102020] = 21
firmware SP = not restored yet
\end{lstlisting}

处理器已经进入固件地址，但返回入口的第一条初始化指令尚未执行。SP 仍表示任务栈的空栈
位置；固件必须先建立自己的 SP，之后才能安全调用发送例程。

\begin{longtable}{|L{0.19\linewidth}|L{0.25\linewidth}|L{0.25\linewidth}|L{0.19\linewidth}|}
\hline
快照 & PC 所属范围 & SP 所属范围 & 当前可执行者 \\ \hline
映像已到齐 & 启动存储 & 固件工作区 & 固件装入器 \\ \hline
任务初始现场 & 任务 RAM & 任务栈 & 求和任务 \\ \hline
返回入口刚进入 & 启动存储 & 任务栈上界 & 固件返回入口 \\ \hline
\end{longtable}

第三行尤其说明 PC 和 SP 不必在同一条指令边界同时“属于同一软件”。受控入口的第一项
职责常常就是从被打断或返回方的栈切换到入口代码自己的栈。当前固定返回协议很简单，但
已经展示了这个状态过渡。

\subsection{错误码怎样对应到最早能够判断的层次}

\begin{longtable}{|L{0.25\linewidth}|L{0.25\linewidth}|L{0.34\linewidth}|}
\hline
错误码 & 最早检测者 & 能证明的事实 \\ \hline
\code{SERIAL\_OVERRUN} & 串行控制器/固件 & 至少一个外部字节未被消费 \\ \hline
\code{HEADER\_INVALID} & 固件头部验证 & 字段不能形成安全布局 \\ \hline
\code{CHECKSUM\_MISMATCH} & 固件载荷验证 & 实收载荷不同于发送端声明 \\ \hline
\code{INVALID\_OPCODE} & 处理器译码 & PC 指向的四字节不是合法指令 \\ \hline
\code{UNMAPPED} & 地址互连 & 一笔合法格式请求没有目标 \\ \hline
\code{ALIGNMENT} & 处理器或互连 & 地址违反当前宽度对齐规则 \\ \hline
\end{longtable}

错误码不能相互替代。校验失败不能说明是哪条指令错误；无效操作码也不能证明串行传输损坏，
因为发送端可能原本就提供了错误代码。诊断应报告最早被可靠观察到的契约破坏，同时保留
PC、地址和装入阶段等上下文。

经过这些计数、快照和不变量检查，本章不再依赖“任务应该可以运行”的直觉。机器的正确
起点、数据移动量、执行结果和失败位置都能由具体状态复算。

\subsection{同一个“完成”对三个观察者含义不同}

求和任务、固件和外部发送端处在不同位置，能够观察的事件也不同。任务在存储指令收到成功
响应后知道结果已经进入 RAM；固件只有在任务返回后才开始读取结果；外部端则要等结果帧
通过串行线到达并完成校验。

\begin{longtable}{|L{0.21\linewidth}|L{0.28\linewidth}|L{0.35\linewidth}|}
\hline
观察者 & 首次能够确认的事件 & 此时仍不能推出 \\ \hline
求和任务 & \code{STORE} 得到 \code{OK} & 固件已经运行、外部端已经收到 \\ \hline
固件返回入口 & 重新读取结果槽得到 21 & 发送线已经传完结果帧 \\ \hline
外部发送端 & 完整结果帧到达且校验通过 & 结果断电后仍会保留 \\ \hline
\end{longtable}

若固件在写入 \code{TXDATA} 后立即复位串行控制器，最后一个字节可能还在移位寄存器中，
外部端收不到完整帧。因而“处理器完成设备寄存器写入”和“设备完成物理发送”也不是同一
时刻。最小接口可以增加 \code{TX\_EMPTY} 状态，表示发送槽和移位寄存器都为空；固件在
进入等待或复位前检查它。

\textbf{硬件模型。} 发送侧除 \code{tx\_ready} 外保存移位寄存器占用状态，最后一位离开
引脚后置 \code{tx\_empty=1}。

\textbf{软件模型。} \code{STATUS} 的 bit3 暴露 \code{TX\_EMPTY}。写
\code{TXDATA} 后该位清零，整字节发送结束后重新置位。

\textbf{抽象。} 软件等待 \code{TX\_EMPTY=1}，可以依赖此前提交的字节已经离开控制器；
这仍不保证远端正确采样或通过帧校验，端到端确认必须由远端协议提供。

\begin{center}
\begin{tikzpicture}[node distance=10mm and 13mm]
\node[stage,minimum width=3.0cm] (store) {RAM 写完成};
\node[stage,minimum width=3.0cm,right=of store] (return) {任务返回固件};
\node[stage,minimum width=3.0cm,right=of return] (queue) {写入 TXDATA};
\node[stage,minimum width=3.0cm,below=13mm of queue] (empty) {TX\_EMPTY 置位};
\node[stage,minimum width=3.0cm,left=of empty] (remote) {远端整帧校验};
\draw[flow] (store)--(return);
\draw[flow] (return)--(queue);
\draw[flow] (queue)--(empty);
\draw[flow] (empty)--(remote);
\end{tikzpicture}
\end{center}
\diagramnote{五个事件顺序发生，分别关闭计算、控制流、设备提交、物理发送和端到端接收边界。}

\subsection{重新运行任务与整机复位有什么不同}

固件收到下一份合法头部后，可以覆盖任务区、重建栈并再次提交入口，而不必让电源复位。
这种重新运行保留启动存储内容和固件初始化后的设备配置，只重建与本次任务相关的状态。

整机复位则把 PC 送回 \code{0x00000000}，清除处理器控制状态、串行协议状态和互连未完成
事务。RAM 字节可能仍在，但不再被信任。两条路径的共同点是：在任务入口提交前都必须重新
验证映像并建立现场。

\begin{longtable}{|L{0.24\linewidth}|L{0.27\linewidth}|L{0.35\linewidth}|}
\hline
状态 & 固件重新运行任务 & 整机复位 \\ \hline
启动存储 & 不变 & 不变 \\ \hline
固件设备配置 & 可以沿用 & 重新初始化 \\ \hline
RAM 旧任务字节 & 按新映像覆盖或清零 & 可能仍在，但必须视为未验证 \\ \hline
处理器 PC/SP & 由固件直接建立新任务现场 & 先回复位入口，再执行固件 \\ \hline
未完成串行事务 & 由协议明确丢弃和重同步 & 控制器复位后丢弃 \\ \hline
\end{longtable}

“重启程序”在当前裸机上不是硬件指令，而是一段固件流程。硬件复位提供确定的最外层恢复
手段，固件重新运行则是较小范围的软件动作。

\subsection{自修改写入说明保护不能由布局表代替}

设错误任务把结果地址 \code{r2} 改成 \code{0x00100004}，即循环中 \code{LOAD} 指令的
位置。最终 \code{STORE} 会把整数 21 的字节 \code{15 00 00 00} 写入代码。任务随后立即
执行位于 \code{0x0010001c} 的 \code{RET}，所以本次仍可能正常回到固件。

固件若不重新装入就再次从 \code{0x00100000} 运行任务，第二条指令字节已经不再是
\code{20 05 00 00}。操作码 \code{0x15} 未定义，处理器进入
\code{INVALID\_OPCODE} 错误入口。第一次运行的结果正确，第二次才暴露代码破坏。

\begin{longtable}{|L{0.20\linewidth}|L{0.28\linewidth}|L{0.36\linewidth}|}
\hline
时刻 & \code{MEM[0x00100004..07]} & 后果 \\ \hline
装入完成 & \code{20 05 00 00} & 译码为数组装入 \\ \hline
错误存储后 & \code{15 00 00 00} & 当前 PC 已越过此处，暂未失败 \\ \hline
再次从入口执行 & 同上 & 第二条取指得到无效操作码 \\ \hline
\end{longtable}

地址布局告诉任务“哪里应当是代码”，RAM 和互连却没有只读规则。要阻止这种写入，机器
需要一种把地址范围与允许操作绑定的强制机制。此处只记录失败，不提前展开地址保护；它会
在后续机器真正需要隔离时成为新的演进问题。

\subsection{本章事实之间的最终依赖关系}

\begin{center}
\begin{tikzpicture}[node distance=9mm and 13mm]
\node[stage,minimum width=3.2cm] (reset) {复位入口可靠};
\node[stage,minimum width=3.2cm,right=of reset] (fw) {固件能够执行};
\node[stage,minimum width=3.2cm,right=of fw] (load) {映像可验证装入};
\node[stage,minimum width=3.2cm,below=13mm of load] (state) {初始现场完整};
\node[stage,minimum width=3.2cm,left=of state] (run) {任务逐指令推进};
\node[stage,minimum width=3.2cm,left=of run] (result) {结果返回并报告};
\draw[flow] (reset)--(fw);
\draw[flow] (fw)--(load);
\draw[flow] (load)--(state);
\draw[flow] (state)--(run);
\draw[flow] (run)--(result);
\end{tikzpicture}
\end{center}
\diagramnote{每个方框都依赖前一个方框已经建立的状态。缺少任意一环，都不能从“字节存在”
直接推出“任务完成”。}

至此，最小机器的器件、连接、软件接口、启动协议、执行 trace 和错误边界均已落到具体状态。
后续章节可以把这台机器作为已知起点，而不必再次把“处理器能够执行一项任务”当作未经证明
的前提。

\subsection{从可观察现象反推连接错误}

具体机器模型还应支持诊断。若上电后没有得到结果，不能只说“机器没跑起来”，而要利用
PC、事务和状态寄存器把故障范围逐步缩小。

\subsection{复位后连第一笔请求都没有}

若示波观察不到 \code{req\_valid}，先检查时钟是否产生边沿、复位是否释放以及 PC 是否
得到规定值。此时还没有理由检查任务映像，因为固件第一条取指尚未发生。

\subsection{请求地址正确，但始终没有响应}

观察到地址 \code{0x00000000} 和读请求，却没有 \code{resp\_valid}，故障位于互连选择、
启动存储连接或目标握手。若互连错误地把该地址送往 RAM，RAM 中的随机字节可能返回成功，
随后表现为无效操作码。区分“无响应”和“错误目标响应”需要同时观察目标选择信号。

\subsection{固件运行，但接收状态永远为零}

若 PC 已在固件轮询循环，\code{STATUS} 读也持续成功，问题不在处理器取指和 RAM。应检查
外部发送端是否产生串行位、控制器采样时序是否匹配以及接收移位寄存器是否能置
\code{RX\_READY}。软件看到的零值只是控制器没有宣布完整字节，不足以判断线路哪一端
失效。

\subsection{任务结果正确，外部端却没有收到}

先读 \code{MEM32[0x00102020]}。若值为 21，计算和 RAM 写入已经完成；再看 PC 是否到达
固件返回入口，区分“任务未返回”和“返回后发送失败”；最后检查
\code{TX\_READY/TX\_EMPTY} 与远端帧校验。诊断顺序沿完成边界向外推进。

\begin{longtable}{|L{0.25\linewidth}|L{0.27\linewidth}|L{0.34\linewidth}|}
\hline
最后确认成功的事实 & 下一项观察 & 尚不应检查的远端事项 \\ \hline
时钟与复位正常 & 第一笔取指请求 & 任务求和结果 \\ \hline
固件取指成功 & 串行状态与接收字节 & 任务返回地址 \\ \hline
映像校验成功 & 初始现场与入口 PC & 外部结果帧 \\ \hline
结果槽为 21 & 返回入口 PC & 传输校验 \\ \hline
\code{TX\_EMPTY=1} & 远端接收与帧校验 & RAM 计算过程 \\ \hline
\end{longtable}

这种从已知成功边界向下一边界推进的方法，与后续操作系统中的系统调用、设备请求和文件
持久化诊断具有相同结构：先确定事实在哪一层成立，再检查下一层所需的新状态。

\subsection{最小机器的最终接线清单}

到本章末再列清单，作用是复核已经推导的连接，而不是用名词代替正文：

\begin{longtable}{|L{0.21\linewidth}|L{0.31\linewidth}|L{0.34\linewidth}|}
\hline
连接 & 承载的信息 & 接错后的直接后果 \\ \hline
时钟到处理器、互连和控制器 & 共同状态更新边沿 & 握手双方不能在同一边界观察状态 \\ \hline
复位到处理器与控制状态 & PC 起点、协议初态 & 旧事务或任意 PC 被当作有效 \\ \hline
处理器请求到互连 & 地址、读写、宽度、写数据 & 无法取指或访问错误目标 \\ \hline
目标响应经互连回处理器 & 状态与读数据 & 指令永久等待或读取错误字节 \\ \hline
互连到启动存储 & 固件读请求 & 复位后没有可靠指令 \\ \hline
互连到 RAM & 任务取指和数据读写 & 任务无法装入或保存运行状态 \\ \hline
互连到串行控制器 & 状态、收发字节 & 外部任务无法进入或结果无法报告 \\ \hline
串行控制器到外部端 & 按时间传送的位 & 软件寄存器与外部信息断开 \\ \hline
\end{longtable}

每一行都能回到前文的一次具体事务。清单中没有中断控制器、计时器、地址转换器、磁盘或
操作系统；这些结构只有在后续问题确实需要时才会加入，并重新说明硬件模型、软件接口和
提供的抽象。


\topic{本章留下的机器}
固定启动程序已经能够接收、验证和装入一项外部程序，并在程序
合作返回后等待下一次传输。但 RAM 中只有一套当前程序布局和一份活动处理器状态。若希望
保留 A 的中间结果，同时让 B 获得处理器，就必须回答暂停现场保存在哪里以及怎样恢复。
